% Part: applied-modal-logics
% Chapter: temporal-logic
% Section: extra-operators

\documentclass[../../../include/open-logic-section]{subfiles}

\begin{document}

\olfileid[fr]{aml}{tl}{ext}

\olsection{Autres opérateurs de la logique temporelle}

Outre les opérateurs unaires du passé et du futur, les logiques
temporelles comprennent parfois des opérateurs binaires $\Since $
et~$\Until$, destinés à exprimer «~depuis~» (\emph{since}) et
«~jusqu'à~» (\emph{until}). Il faut alors ajouter $\Since $ et~$\Until$
au langage de la logique temporelle, ainsi que la clause suivante à
la définition d'une !!{formula} temporelle :

\begin{itemize}
\item[] Si $!A$ et $!B$ sont des !!{formula}s, alors $(\Since  !A !B)$
  et $(\Until !A !B)$ sont toutes deux des !!{formula}s.
\end{itemize}

La sémantique de ces opérateurs est définie comme suit :

\begin{defn}\ollabel{defn:since-until}
  \emph{Vérité d'!!a{formula}~$!A$ à l'instant~$t$} dans un modèle~$\mModel M$ :
  \begin{enumerate}
  \item\ollabel{defn:sub:mmodels-since} \indcase{!A}{\Since  !B
    !C}{$\mSat{M}{\indfrm}[t]$ si et seulement s'il existe $t' \in T$
    tel que $t' \prec t$, avec $\mSat{M}{!B}[t']$ et, pour tout $s$
    tel que $t' \prec s \prec t$, $\mSat{M}{!C}[s]$}
  \item\ollabel{defn:sub:mmodels-until} \indcase{!A}{\Until !B
    !C}{$\mSat{M}{\indfrm}[t]$ si et seulement s'il existe $t' \in T$
    tel que $t \prec t'$, avec $\mSat{M}{!B}[t']$ et, pour tout $s$
    tel que $t \prec s \prec t'$, $\mSat{M}{!C}[s]$}
  \end{enumerate}
\end{defn}

La lecture intuitive de $\Since  !B !C$ est «~Depuis que $!B$ a été
le cas, $!C$~a été le cas.~» Celle de $\Until !B !C$ est «~Jusqu'à
ce que $!B$ soit le cas, $!C$~sera le cas.~»

\end{document}
